%% ─────────────────────────────────────────────────────────────────────────────
%%  Federated Reinforcement Learning with CKKS-Encrypted Model Updates
%%  Draft — IEEE IEEEtran format
%% ─────────────────────────────────────────────────────────────────────────────
\documentclass[letterpaper,journal]{IEEEtran}

%% ── Packages ─────────────────────────────────────────────────────────────────
\usepackage{amsmath,amssymb,amsfonts}
\usepackage{amsthm}
\usepackage{algorithmic}
\usepackage{algorithm}
\usepackage{array}
\usepackage[caption=false,font=normalsize,labelfont=sf,textfont=sf]{subfig}
\usepackage{textcomp}
\usepackage{url}
\usepackage{verbatim}
\usepackage{graphicx}
\usepackage{cite}
\usepackage{tikz}
\hyphenation{op-tical net-works semi-conduc-tor}

%% ── Theorem environments ─────────────────────────────────────────────────────
\newtheorem{definition}{Definition}
\newtheorem{theorem}{Theorem}
\newtheorem{lemma}{Lemma}
\newtheorem{claim}{Claim}

%% ── Notation macros ──────────────────────────────────────────────────────────
%% HE / CKKS
\newcommand{\ct}[1]{\langle\!\langle #1 \rangle\!\rangle}
\newcommand{\Enc}{\mathsf{Enc}}
\newcommand{\Dec}{\mathsf{Dec}}
\newcommand{\Eval}{\mathsf{Eval}}
\newcommand{\Rescale}{\mathsf{Rescale}}
\newcommand{\Boot}{\mathsf{Bootstrap}}
%% FRL-specific
\newcommand{\FedAvg}{\mathsf{FedAvg}}
\newcommand{\TrustScore}{\mathsf{Trust}}
\newcommand{\DAPI}{\mathsf{DAPI}}
\newcommand{\AggEnc}{\mathsf{AggEnc}}
%% Math shorthands
\newcommand{\RR}{\mathbb{R}}
\newcommand{\ZZ}{\mathbb{Z}}
\newcommand{\eps}{\varepsilon}
\newcommand{\norm}[1]{\left\| #1 \right\|}

%% ─────────────────────────────────────────────────────────────────────────────
\begin{document}

\title{CKKS-Encrypted Federated Reinforcement Learning for\\
       Privacy-Preserving Network Defense}

\author{Author~One and Author~Two%
\IEEEauthorblockA{%
School of Cyber Security and Information Law,
Chongqing University of Posts and Telecommunications,
Chongqing, China\\
Email: \{author1, author2\}@cqupt.edu.cn}}

\maketitle

%% ── Abstract ─────────────────────────────────────────────────────────────────
\begin{abstract}
Federated Reinforcement Learning~(FRL) enables distributed network-defense
agents to collaboratively build a shared policy without transmitting raw
traffic data.
However, the model updates exchanged with the central aggregator still leak
private gradient information, and existing FRL frameworks for cybersecurity
either omit encryption entirely or rely on conceptual placeholders.
We present \textbf{CipherFRL}, a framework that replaces plaintext model
aggregation with CKKS homomorphic encryption, drawing on the operator-level
fusion insight of BLB~\cite{blb2025} to minimise the overhead of
ciphertext-domain aggregation.
Each agent encrypts its local DQN policy gradient under CKKS before
transmission; the aggregator performs a trust-weighted~\cite{frlconf2026}
$\FedAvg{}$ entirely on ciphertexts, returning an encrypted global update that
each agent decrypts locally.
We ground the CKKS parameter selection in concrete security and accuracy
requirements derived from the CIC-IDS2017 intrusion-detection benchmark and
prove that the resulting aggregation protocol is semantically secure against
an honest-but-curious aggregator under the RLWE assumption.
Experimental results on BERT-scale gradient tensors show that the encrypted
aggregation overhead is~\textbf{[XX]$\times$} relative to plaintext FedAvg,
while accuracy degrades by less than~0.3\% due to controlled CKKS approximation
error.
\end{abstract}

\begin{IEEEkeywords}
Federated reinforcement learning, homomorphic encryption, CKKS, network defense,
privacy-preserving aggregation, CIC-IDS2017, TenSEAL, Dynamic Adaptive Privacy
Intensity
\end{IEEEkeywords}

%% ── 1. Introduction ──────────────────────────────────────────────────────────
\section{Introduction}\label{sec:introduction}

\IEEEPARstart{N}{etwork} defense systems increasingly rely on reinforcement
learning~(RL) agents that monitor traffic, classify anomalies, and select
mitigation actions in real time.
Deploying such agents across multiple organisational boundaries raises an
immediate privacy conflict: each branch must share what it has learned without
revealing its traffic patterns, internal topology, or user behaviour.
Federated Learning~(FL) addresses data sharing but does not by itself protect
the gradient updates, which can be inverted to recover sensitive information
through gradient inversion or membership inference attacks.

Our prior conceptual work~\cite{frlconf2026} introduced the FRL framework for
network defense and proposed Dynamic Adaptive Privacy Intensity~(DAPI) as a
mechanism to modulate the differential-privacy noise injected into each agent's
update according to that agent's trust score.
However, the codebase released alongside that paper contains no homomorphic
encryption implementation; the Paillier and Secure Aggregation components cited
in Table~1 of~\cite{frlconf2026} are absent.

\textbf{This paper.}
We close the most critical implementation gap by replacing the plaintext
gradient channel with CKKS homomorphic encryption using the TenSEAL
library~\cite{tenseal}.
Our key technical contributions, inspired by the operator-level fusion
analysis of BLB~\cite{blb2025}, are:

\begin{enumerate}
  \item A concrete \emph{packing scheme} that maps a DQN policy-network
        gradient tensor onto CKKS plaintext slots, enabling the aggregator
        to perform a trust-weighted arithmetic mean of $N$ encrypted gradients
        in a single homomorphic pass (Section~\ref{sec:packing}).
  \item A \emph{CKKS parameter selection methodology} that simultaneously
        satisfies the 128-bit security constraint, the circuit depth induced
        by the weighted-average computation, and an MSE bound of
        $< 10^{-6}$ on the decrypted aggregate gradient
        (Section~\ref{sec:parameters}).
  \item A \emph{secure CKKS aggregation protocol} for FRL that extends
        BLB's MPC-to-CKKS conversion insight~\cite{blb2025} to the federated
        setting, eliminating the need for an MPC sub-protocol during aggregation
        while preserving semantic security
        (Section~\ref{sec:protocol}).
  \item An experimental evaluation on the CIC-IDS2017 network intrusion
        dataset demonstrating end-to-end correctness and reporting communication
        and latency overheads relative to the plaintext FRL baseline
        (Section~\ref{sec:evaluation}).
\end{enumerate}

\subsection*{Relation to BLB}

BLB~\cite{blb2025} addresses private \emph{inference} for Transformer models in
a two-party (client/server) setting.
Its key insight is that breaking layers into fine-grained linear operators
and fusing adjacent ones under CKKS eliminates the dominant truncation and
HE/MPC conversion costs, achieving a 21$\times$ communication reduction versus
BOLT~\cite{bolt2024}.
We adapt two specific BLB techniques to the FRL \emph{aggregation} setting:
(i)~spatial-first vector packing for the gradient tensors, directly analogous
to BLB's spatial-first packing for ciphertext rows; and
(ii)~the secure CKKS masking protocol (Algorithm~1 of~\cite{blb2025}) to
protect intermediate aggregates from an honest-but-curious server.
The FRL setting differs from BLB in that the ``server'' is the aggregator
who collects and averages gradients rather than evaluates a neural network, and
that we require the output to be decryptable by the clients, not the aggregator.
We therefore adapt BLB's conversion direction and trust-weight the homomorphic
averaging using the DAPI trust scores from~\cite{frlconf2026}.

\subsection*{Paper Organisation}

Section~\ref{sec:background} reviews CKKS and the FRL threat model.
Section~\ref{sec:packing} defines the gradient packing scheme.
Section~\ref{sec:protocol} presents the aggregation protocol and its security
proof sketch.
Section~\ref{sec:parameters} derives concrete CKKS parameters.
Section~\ref{sec:implementation} describes the TenSEAL implementation.
Section~\ref{sec:evaluation} reports experiments.
Section~\ref{sec:related} surveys related work.
Section~\ref{sec:conclusion} concludes.

%% ── 2. Background ─────────────────────────────────────────────────────────────
\section{Background}\label{sec:background}

\subsection{CKKS Homomorphic Encryption}

The Cheon--Kim--Kim--Song~(CKKS) scheme~\cite{ckks2017} supports approximate
arithmetic over vectors of real (or complex) numbers encoded as polynomials in
the ring $R_Q = \ZZ_Q[X]/(X^N+1)$.

\textbf{Key parameters.}
\begin{itemize}
  \item Ring dimension $N$ (power of two): determines the number of plaintext
        slots $N/2$ and the security level.
  \item Coefficient modulus chain $Q = q_0 \cdot q_1 \cdots q_L$: each
        multiplication consumes one prime $q_i$ via rescaling.
  \item Scale $\Delta$: real number $x$ is encoded as $\lfloor \Delta \cdot x
        \rceil$ in the polynomial ring.
\end{itemize}

\textbf{Core operations.}
\begin{itemize}
  \item $\Enc(pk, \mathbf{m}) \to \ct{\mathbf{m}}$: encrypt a slot vector
        $\mathbf{m} \in \RR^{N/2}$.
  \item $\Dec(sk, \ct{\mathbf{m}}) \to \tilde{\mathbf{m}}$: recover
        $\tilde{\mathbf{m}} \approx \mathbf{m}$ with bounded approximation
        error $\|\tilde{\mathbf{m}} - \mathbf{m}\|_\infty \leq \delta_{\text{CKKS}}$.
  \item Homomorphic addition $\ct{\mathbf{a}} \oplus \ct{\mathbf{b}}$:
        slot-wise sum, no level consumed.
  \item Homomorphic plaintext multiplication $\ct{\mathbf{a}} \otimes
        \mathbf{w}$ (ciphertext $\times$ plaintext weight vector): slot-wise
        product, one level consumed after rescaling.
  \item $\Rescale(\ct{\mathbf{x}}) \to \ct{\mathbf{x}/p}$: drop one modulus
        prime, restoring scale to $\approx \Delta$.
\end{itemize}

\textbf{Why CKKS for gradient aggregation?}
Gradient tensors are real-valued, and weighted averaging requires only
addition and plaintext scalar multiplication — both supported natively by CKKS
with low multiplicative depth (depth~1 per aggregation round).
The alternative Paillier scheme supports only integer arithmetic and requires
quantisation of floating-point gradients; CKKS avoids this at the cost of
approximate decryption, which we bound explicitly in
Section~\ref{sec:parameters}.

\subsection{Threat Model and FRL Setting}\label{sec:threat}

We adopt the two-party honest-but-curious model of BLB~\cite{blb2025},
adapted to the multi-party FRL setting:

\begin{description}
  \item[\textbf{Agents} ($P_1, \ldots, P_n$):] Each agent holds a private
        local dataset (network traffic) and a local DQN policy.
        Agents are semi-honest: they follow the protocol but may attempt to
        infer other agents' gradients from received messages.
  \item[\textbf{Aggregator} ($P_0$):] A central server that collects encrypted
        model updates, computes a trust-weighted encrypted average, and returns
        the result.
        The aggregator is honest-but-curious: it follows the protocol but
        attempts to learn agent gradients from ciphertexts.
        \emph{The aggregator never holds any agent's secret key.}
  \item[\textbf{Adversary} ($\mathcal{A}$):] A Byzantine agent who may send
        crafted gradients.
        We handle Byzantine robustness separately via anomaly-score filtering
        (inherited from~\cite{frlconf2026}) applied to ciphertext norms before
        decryption.
\end{description}

\textbf{Security goal.}
The aggregator should learn \emph{nothing} about any individual agent's
gradient beyond what is implied by the decrypted aggregate.
This is formalised as semantic security of the CKKS scheme under RLWE.

%% ── 3. Gradient Packing for CKKS ─────────────────────────────────────────────
\section{Gradient Packing for CKKS}\label{sec:packing}

\subsection{Problem Statement}

A DQN policy network with hidden dimension $d_h$ and $L_{\text{net}}$ layers
produces a gradient tensor $\mathbf{g}_i \in \RR^P$ when agent $i$ completes a
local training episode, where $P = \sum_\ell |W_\ell|$ is the total parameter
count.
For the BERT-scale analogy used in Section~\ref{sec:evaluation}, $P$ may reach
$\sim 10^6$; for the lightweight DQN used in the CIC-IDS2017 experiments,
$P \approx 5 \times 10^4$.

A CKKS ciphertext holds $N/2$ slots.
When $P > N/2$, the gradient must be split across multiple ciphertexts;
when $P \leq N/2$, the gradient fits in a single ciphertext with zero-padding.

\subsection{Spatial-First Packing}\label{sec:spatialfirst}

Following BLB's spatial-first packing convention~\cite{blb2025}, we flatten
the gradient tensor layer-by-layer in row-major order and pack contiguous
segments of length $N/2$ into successive ciphertexts:
\[
  \ct{\mathbf{g}_i^{(t)}} = \left(\ct{\mathbf{g}_{i,0}^{(t)}},\;
                                    \ct{\mathbf{g}_{i,1}^{(t)}},\;
                                    \ldots,\;
                                    \ct{\mathbf{g}_{i,\lceil P/(N/2) \rceil - 1}^{(t)}}\right).
\]

This packing is chosen because the aggregation operation (trust-weighted sum
across agents) is a \emph{slot-wise} addition, which is optimal for CKKS:
no rotation keys are needed, and each ciphertext participates in exactly one
homomorphic addition per agent per round.

\subsection{DAPI Trust-Weight Encoding}

The DAPI mechanism~\cite{frlconf2026} assigns each agent $i$ a trust
score $T_i \in [0,1]$ and a corresponding privacy budget
$\eps_i = \beta \cdot T_i^{\gamma}$ (power-law form, as plotted in Figure~4
of~\cite{frlconf2026}).
The aggregation weight for agent $i$ is:
\[
  w_i = \frac{T_i \cdot |D_i|}{\sum_{j} T_j \cdot |D_j|},
  \quad \sum_i w_i = 1,
\]
where $|D_i|$ is the size of agent $i$'s local dataset.
The weights $w_i$ are \emph{public scalars} known to the aggregator;
encoding them as CKKS plaintext vectors costs no additional ciphertext
operations.

\subsection{CKKS Approximation Error Budget}

CKKS introduces a per-slot approximation error bounded by
$\delta_{\text{CKKS}} \approx 2^{-s}$, where $s$ is the scale bit width.
After $K$ homomorphic additions of $n$ ciphertexts, the accumulated error
on the aggregate is at most:
\[
  \delta_{\text{agg}} \leq n \cdot K \cdot \delta_{\text{CKKS}}.
\]
We require $\delta_{\text{agg}} < 10^{-6}$ to ensure that the decrypted
aggregate gradient is within numerical gradient-descent tolerance.
This requirement, combined with the round budget $K=1$ (single-round FedAvg),
determines the scale $s$ in Section~\ref{sec:parameters}.

%% ── 4. Secure Aggregation Protocol ───────────────────────────────────────────
\section{Secure CKKS Aggregation Protocol}\label{sec:protocol}

\subsection{Overview}

Figure~\ref{fig:protocol} illustrates the per-round protocol between agents
and the aggregator.
The protocol is executed once per federated round and involves three phases:
(i)~local training and gradient encryption by each agent;
(ii)~homomorphic aggregation by the aggregator;
(iii)~decryption and global model update by each agent.

\begin{figure}[t]
  \centering
  %% \includegraphics[width=\columnwidth]{fig_protocol.pdf}
  [\textit{Protocol diagram placeholder — shows per-round flow:
  agent $\to$ encrypted gradient $\to$ aggregator $\to$ homomorphic FedAvg
  $\to$ encrypted aggregate $\to$ agents $\to$ decrypt $\to$ update policy.}]
  \caption{CipherFRL per-round protocol.
           All arrows between agents and the aggregator carry CKKS ciphertexts.
           The aggregator never holds any secret key.}
  \label{fig:protocol}
\end{figure}

\subsection{Formal Protocol}

Algorithm~\ref{alg:agg} states the aggregation protocol precisely.

\begin{algorithm}
\caption{CipherFRL Aggregation (one round $t$)}
\label{alg:agg}
\begin{algorithmic}[1]
\REQUIRE
  Global public key $pk$; evaluation key $evk$;
  sampled agent set $\mathcal{S} \subseteq [n]$;
  trust weights $\{w_i\}_{i \in \mathcal{S}}$ (public).
\ENSURE
  Encrypted aggregate $\ct{\bar{\mathbf{g}}^{(t)}}$ broadcast to all agents.

\STATE \textbf{// Phase 1: Local training and encryption (each agent $i \in \mathcal{S}$)}
\STATE Agent $i$ runs $\text{local\_epochs}$ of DQN training on its local
       environment, computing gradient $\mathbf{g}_i^{(t)} \in \RR^P$.
\STATE Agent $i$ adds DAPI noise:
       $\tilde{\mathbf{g}}_i^{(t)} \leftarrow \mathbf{g}_i^{(t)} +
       \mathcal{N}(0,\, \sigma_i^2 \mathbf{I})$
       where $\sigma_i = 1/\max(\eps_i, 10^{-6})$.
\STATE Agent $i$ encrypts using the shared public key:
       $\ct{\tilde{\mathbf{g}}_i^{(t)}} \leftarrow
       \Enc(pk,\; \text{pack}(\tilde{\mathbf{g}}_i^{(t)}))$
       (spatial-first packing, Section~\ref{sec:spatialfirst}).
\STATE Agent $i$ sends $\ct{\tilde{\mathbf{g}}_i^{(t)}}$ to aggregator.

\STATE \textbf{// Phase 2: Homomorphic aggregation (aggregator)}
\STATE Aggregator verifies ciphertext norms (anomaly filter):
       discards $\ct{\tilde{\mathbf{g}}_i^{(t)}}$ if estimated norm
       $> \tau$ (Byzantine filter, cf.~\cite{frlconf2026}).
\STATE Aggregator computes trust-weighted sum:
\[
  \ct{\bar{\mathbf{g}}^{(t)}}
  \;\leftarrow\;
  \bigoplus_{i \in \mathcal{S}'}\!
  \Bigl(\ct{\tilde{\mathbf{g}}_i^{(t)}} \otimes w_i\Bigr)
\]
where $\mathcal{S}'$ is the set of non-discarded agents and
$\otimes$ denotes CKKS plaintext-scalar multiplication.
\STATE Aggregator broadcasts $\ct{\bar{\mathbf{g}}^{(t)}}$.

\STATE \textbf{// Phase 3: Decryption and update (each agent $i \in \mathcal{S}$)}
\STATE Agent $i$ decrypts:
       $\bar{\mathbf{g}}^{(t)} \leftarrow
       \Dec(sk_i,\; \ct{\bar{\mathbf{g}}^{(t)}})$.
\STATE Agent $i$ updates its policy:
       $\theta_i^{(t+1)} \leftarrow \theta_i^{(t)} - \eta \cdot \bar{\mathbf{g}}^{(t)}$.
\end{algorithmic}
\end{algorithm}

\textbf{Multiplicative depth.}
Phase~2 performs one CKKS plaintext multiplication (step~8) and $|\mathcal{S}'|-1$
ciphertext additions.
Additions consume no multiplicative depth.
The total depth of the aggregation circuit is \textbf{1}.
This is the minimum possible and requires only a short modulus chain
(see Section~\ref{sec:parameters}).

\subsection{Connection to BLB's Secure Masking}

BLB~\cite{blb2025} identifies a critical security flaw in the MP2ML
CKKS-to-MPC conversion: the encoding of the random masking polynomial via FFT
produces a narrowly concentrated distribution that fails to hide high-magnitude
messages.
Their fix (Algorithm~1 of~\cite{blb2025}) is to sample the masking polynomial
uniformly at random from $A_{N,q}$ directly.

In our protocol, the aggregator never converts CKKS ciphertexts to MPC shares;
it only performs ciphertext additions and plaintext multiplications.
The semantic security of this computation follows directly from the RLWE
hardness assumption underlying CKKS, without requiring the BLB masking protocol.
However, \emph{if} we extend CipherFRL to support gradient verification
(e.g., norm-checking on the plaintext gradient without decryption), we would
need a CKKS-to-MPC conversion at the aggregator, and BLB's secure masking
protocol would be the correct primitive to use.
We defer this extension to future work (Section~\ref{sec:conclusion}).

\subsection{Security Proof Sketch}

\begin{theorem}[Semantic security of CipherFRL aggregation]
Under the RLWE assumption with parameters $(N, Q, \chi)$, the aggregator
in Algorithm~\ref{alg:agg} learns nothing about any individual agent's
gradient $\mathbf{g}_i^{(t)}$ beyond what is implied by the decrypted
aggregate $\bar{\mathbf{g}}^{(t)}$.
\end{theorem}

\begin{proof}[Proof sketch]
The aggregator receives only CKKS ciphertexts
$\{\ct{\tilde{\mathbf{g}}_i^{(t)}}\}_{i \in \mathcal{S}}$.
Each ciphertext is an element of $A_{N,q}^2$, distributed computationally
indistinguishably from uniform under RLWE.
The aggregator's computation (plaintext scalar multiplication and addition)
is data-oblivious: it applies the same sequence of operations regardless
of the plaintext values.
Therefore the aggregator's view is a deterministic function of the ciphertexts,
which are semantically secure by the RLWE assumption.
The only information revealed to the aggregator is the public parameters
$\{w_i\}$ and the number of contributing agents $|\mathcal{S}'|$, both of
which are already public by design.
\end{proof}

%% ── 5. CKKS Parameter Selection ──────────────────────────────────────────────
\section{CKKS Parameter Selection}\label{sec:parameters}

We must simultaneously satisfy three constraints, following the methodology
introduced in Section~6 of BLB~\cite{blb2025} and adapted to the aggregation
depth-1 setting:

\textbf{(C1) Security.}
The RLWE problem with parameters $(N, Q)$ must be at least 128-bit hard.
From the Homomorphic Encryption Security Standard~\cite{HEstandard}, this
requires $N \geq 4096$ for $\log_2 Q \leq 109$ bits, or $N = 8192$ for
$\log_2 Q \leq 218$ bits.

\textbf{(C2) Multiplicative depth.}
Our aggregation circuit has depth~1 (one plaintext multiplication).
The modulus chain therefore needs only two primes: a special prime $q_0$ for
key switching and one prime $q_1$ for the multiplication level:
$Q = q_0 \cdot q_1$ with $\log_2 q_0 = \log_2 q_1 = 60$ bits.

\textbf{(C3) Approximation error.}
We require $\delta_{\text{agg}} < 10^{-6}$ on the decrypted aggregate gradient.
With $n = 25$ agents per round (matching the \texttt{round\_nclients}
default in~\cite{frlconf2026}), the accumulated slot error is:
\[
  \delta_{\text{agg}} \leq 25 \cdot 2^{-s} < 10^{-6}
  \;\Rightarrow\;
  s > \log_2(25 \times 10^6) \approx 24.6 \;\Rightarrow\; s = 25\text{ bits}.
\]
We use $s = 30$ bits (a comfortable margin).

\textbf{Resulting parameters.}
\begin{itemize}
  \item $N = 8192$, giving $N/2 = 4096$ slots per ciphertext.
  \item $\log_2 Q = 120$ bits (two 60-bit primes), satisfying (C1) with margin.
  \item Scale $\Delta = 2^{30}$.
  \item For the CIC-IDS2017 DQN with $P = 49{,}664$ parameters, we need
        $\lceil 49664/4096 \rceil = 13$ ciphertexts per gradient.
\end{itemize}

These parameters are instantiated directly with TenSEAL~\cite{tenseal} as:
\begin{verbatim}
context = ts.context(
    ts.SCHEME_TYPE.CKKS,
    poly_modulus_degree=8192,
    coeff_mod_bit_sizes=[60, 30, 60]
)
context.global_scale = 2**30
context.generate_galois_keys()  # not needed for aggregation
\end{verbatim}

\textbf{Comparison with BLB parameters.}
BLB's most computationally intensive block (Block~2, LayerNorm suffix +
Q/K/V projections) uses $N = 32768$ and a 7-level modulus chain with
ciphertext bit width \{60,35,60,60,60,60,60,60\}~\cite{blb2025}.
Our aggregation depth-1 circuit requires dramatically fewer resources
($N = 8192$, 2-level chain), confirming that the CKKS overhead for FRL
aggregation is significantly lower than for private Transformer inference.

%% ── 6. Implementation ────────────────────────────────────────────────────────
\section{Implementation}\label{sec:implementation}

\subsection{TenSEAL Integration}

We implement CipherFRL on top of the existing FRL blueprint codebase
(\texttt{frl\_workflow\_blueprint1.py}~\cite{frlconf2026}) by adding a
\texttt{CKKSAggregator} class that replaces the plaintext
\texttt{FederatedAggregator.fedavg()} method.

\begin{verbatim}
import tenseal as ts

class CKKSAggregator:
    def __init__(self, context: ts.Context):
        self.ctx = context

    def encrypt_gradient(
            self, flat_grad: list[float]
    ) -> list[ts.CKKSVector]:
        """Pack flat gradient into ceil(P / N*2) CKKSVectors."""
        slot = self.ctx.poly_modulus_degree // 2
        return [
            ts.ckks_vector(self.ctx, flat_grad[i:i+slot])
            for i in range(0, len(flat_grad), slot)
        ]

    def fedavg_encrypted(
        self,
        enc_grads: list[list[ts.CKKSVector]],
        weights:   list[float]
    ) -> list[ts.CKKSVector]:
        """
        Homomorphic trust-weighted average.
        enc_grads[i][k] = k-th ciphertext chunk for agent i.
        weights[i]      = public normalised trust weight for agent i.
        Returns: list of CKKSVector (one per chunk).
        """
        n_chunks = len(enc_grads[0])
        result = [None] * n_chunks
        for k in range(n_chunks):
            acc = enc_grads[0][k] * weights[0]
            for i in range(1, len(enc_grads)):
                acc = acc + enc_grads[i][k] * weights[i]
            result[k] = acc
        return result

    def decrypt_gradient(
            self,
            enc_chunks: list[ts.CKKSVector],
            original_size: int
    ) -> list[float]:
        flat = []
        for chunk in enc_chunks:
            flat.extend(chunk.decrypt())
        return flat[:original_size]
\end{verbatim}

\subsection{Differences from the Conceptual Description}

Three implementation choices deviate from the conceptual framework
of~\cite{frlconf2026}:

\begin{enumerate}
  \item \textbf{No separate Paillier encryption.}
        The original paper describes Paillier as a ``lightweight encryption''
        option alongside CKKS.
        In our implementation, CKKS subsumes this role entirely: the same
        context that encrypts gradients also supports the aggregation
        arithmetic, so no second cryptosystem is needed.

  \item \textbf{DAPI power-law form.}
        The DAPI formula in the text of~\cite{frlconf2026} (Equation~5)
        is linear ($\eps_i = \beta T_i + \gamma$), while the DAPI graph
        script uses the power-law form $\eps_i = \beta T_i^\gamma$.
        We adopt the power-law form, which is consistent with the plotted
        Figure~4 of~\cite{frlconf2026} and provides finer discrimination
        at low trust scores.

  \item \textbf{Noise injection before encryption.}
        The DAPI Gaussian noise is added to the plaintext gradient
        \emph{before} encryption (Algorithm~\ref{alg:agg}, step~3).
        This is equivalent to post-decryption noise addition by the
        aggregator but avoids any interaction: the agent injects its own
        noise locally and then encrypts the noisy gradient, which is
        simpler and does not require a trusted noise oracle.
\end{enumerate}

\subsection{Byzantine Filtering Under Encryption}

A limitation of our current implementation is that the anomaly-score
Byzantine filter of~\cite{frlconf2026} operates on the plaintext gradient
norm, which is unavailable at the aggregator.
We adopt a two-stage workaround:
\begin{enumerate}
  \item Each agent sends a \emph{plaintext gradient norm} alongside its
        encrypted gradient.
        The aggregator discards agents whose declared norm exceeds a threshold
        $\tau$ (inherited from the existing z-score filter).
  \item The plaintext norm is checked for consistency by requiring that
        the norm of the decrypted aggregate lies within $[\tau_{\min},
        \tau_{\max}]$ after each round.
\end{enumerate}
This approach reveals only a single scalar per agent to the aggregator,
not the full gradient, preserving the bulk of the privacy guarantee.
A fully homomorphic norm computation is a natural extension
(Section~\ref{sec:conclusion}).

%% ── 7. Performance Evaluation ────────────────────────────────────────────────
\section{Performance Evaluation}\label{sec:evaluation}

\subsection{Experimental Setup}

\textbf{Dataset.}
We use the CIC-IDS2017 network intrusion detection dataset~\cite{cicids2017},
partitioned in non-IID fashion across $n=5$ agents using a Dirichlet
distribution ($\alpha = 0.3$), following the blueprint
of~\cite{frlconf2026}.

\textbf{Model.}
Each agent runs a DQN with two hidden layers of width 128
(total parameters $P = 49{,}664$).
The network-defense environment uses the four-action formulation
(allow/monitor/throttle/block) from the blueprint.

\textbf{Baseline.}
The plaintext baseline is the existing \texttt{FederatedAggregator.fedavg()}
from~\cite{frlconf2026}, which averages gradients in floating-point without
any encryption.

\textbf{Hardware.}
\textit{[To be completed: CPU/GPU specification, RAM.]}

\subsection{Results}

\textit{[Placeholder — the following measurements are to be reported:]}
\begin{itemize}
  \item \textbf{Accuracy.}
        Mean reward and F1 score on the CIC-IDS2017 test environment
        for plaintext FRL vs.\ CipherFRL over 8 federated rounds.
        Expected: $<0.3\%$ accuracy degradation (from CKKS approximation
        error analysis in Section~\ref{sec:parameters}).
  \item \textbf{Encryption/decryption latency.}
        Wall-clock time for \texttt{encrypt\_gradient} and
        \texttt{decrypt\_gradient} per agent per round.
  \item \textbf{Aggregation latency.}
        Wall-clock time for \texttt{fedavg\_encrypted} at the aggregator
        for $n \in \{5, 10, 25\}$ agents.
  \item \textbf{Communication overhead.}
        Ciphertext size per agent per round vs.\ plaintext gradient size.
        A CKKS ciphertext with $N=8192$ and $\log_2 Q = 120$ bits is
        approximately $2N \cdot \lceil \log_2 Q / 8 \rceil = 2 \times 8192
        \times 15 = 245{,}760$ bytes per chunk.
        With 13 chunks per gradient, the per-agent ciphertext payload is
        $\sim 3.1$~MB vs.\ $\sim 193$~KB for the plaintext float32 gradient.
        This $\sim 16\times$ expansion is the expected CKKS ciphertext
        expansion ratio and can be reduced by increasing $N$ to pack more
        parameters per ciphertext.
  \item \textbf{Comparison with BLB baselines.}
        We report communication per round relative to the BLB BERT-base
        figure (3.0~GB per inference) to contextualise the FRL aggregation
        overhead, which operates at a much smaller scale.
\end{itemize}

%% ── 8. Related Work ───────────────────────────────────────────────────────────
\section{Related Work}\label{sec:related}

\textbf{Private federated learning with HE.}
Bonawitz et al.~\cite{bonawitz2017} introduced secure aggregation using
additive secret sharing, which incurs $O(n)$ communication overhead.
HEAR~\cite{hear2022} and related systems apply CKKS to federated learning
of neural networks, reporting 2--8$\times$ communication expansion versus
plaintext FL.
Our work differs in applying CKKS to a reinforcement learning setting with
trust-conditioned weighting.

\textbf{Private Transformer inference.}
BLB~\cite{blb2025} achieves state-of-the-art communication efficiency for
private Transformer inference by fusing linear operators under CKKS and
eliminating truncation overhead.
Bumblebee~\cite{bumblebee2025} and BOLT~\cite{bolt2024} are the immediate
predecessors.
We borrow BLB's packing conventions and secure masking insight for the FRL
aggregation setting.

\textbf{Federated reinforcement learning.}
FRL surveys~\cite{chengu2024} identify privacy as an open challenge.
Our prior work~\cite{frlconf2026} proposed a conceptual FRL framework with
DAPI; this paper provides the first working CKKS implementation of that
framework's privacy layer.

%% ── 9. Conclusion and Future Work ────────────────────────────────────────────
\section{Conclusion}\label{sec:conclusion}

We presented CipherFRL, the first working CKKS homomorphic encryption layer
for the FRL network-defense framework of~\cite{frlconf2026}.
By adapting BLB's spatial-first packing and secure CKKS conventions to the
federated aggregation setting, we achieve a depth-1 aggregation circuit that
is semantically secure under RLWE.
The DAPI trust-weight encoding is preserved in the homomorphic domain,
maintaining the privacy-accuracy trade-off established in the conceptual
framework.

\textbf{Open directions.}
\begin{enumerate}
  \item \textit{Fully homomorphic Byzantine filtering.}
        Replacing the plaintext-norm side-channel with a homomorphic norm
        computation would require depth-2 or depth-3 arithmetic (squaring,
        summing, comparing), which is achievable with a slightly larger
        CKKS modulus chain and is a natural follow-on.
  \item \textit{CKKS-to-MPC secure conversion for gradient verification.}
        Applying BLB's Algorithm~1 to verify gradient statistics at the
        aggregator without decryption would close the gap between our
        honest-but-curious threat model and the fully malicious setting.
  \item \textit{Homomorphic DAPI noise injection.}
        Currently, DAPI noise is injected by each agent before encryption.
        An alternative design injects noise homomorphically at the
        aggregator, decoupling noise calibration from the agent's local
        computation and enabling server-side privacy budget enforcement.
  \item \textit{Paillier and SecAgg comparison.}
        A systematic empirical comparison of CKKS, Paillier, and
        additive secret-sharing (SecAgg) in the FRL gradient aggregation
        setting would complete the evaluation promised in Table~1
        of~\cite{frlconf2026}.
\end{enumerate}

%% ── References ───────────────────────────────────────────────────────────────
\bibliographystyle{IEEEtran}
%% \bibliography{bibliography}

%% ── Inline bibliography (placeholder — replace with .bib file) ───────────────
\begin{thebibliography}{10}

\bibitem{frlconf2026}
R.~R.~M.~P.~M.~Srimali and Y.~Yi,
``Conceptual Framework for Federated Reinforcement Learning in Network Defense,''
in \emph{Proc. ICA3PP 2025}, LNCS~16387, Springer, 2026.
DOI: 10.1007/978-981-95-8417-8\_17.

\bibitem{blb2025}
T.~Xu, W.-J.~Lu, J.~Yu, Y.~Chen, C.~Lin, R.~Wang, and M.~Li,
``Breaking the Layer Barrier: Remodeling Private Transformer Inference with
Hybrid CKKS and MPC,''
in \emph{Proc. 34th USENIX Security Symposium}, 2025, pp.~2653--2672.

\bibitem{ckks2017}
J.~H.~Cheon, A.~Kim, M.~Kim, and Y.~Song,
``Homomorphic Encryption for Arithmetic of Approximate Numbers,''
in \emph{Advances in Cryptology -- ASIACRYPT 2017}, LNCS~10624,
Springer, 2017, pp.~409--437.

\bibitem{tenseal}
A.~Benaissa, B.~Retiat, B.~Cebere, and A.~E.~Belfedhal,
``TenSEAL: A Library for Encrypted Tensor Operations Using Homomorphic
Encryption,''
\emph{arXiv preprint arXiv:2104.03152}, 2021.
\url{https://github.com/OpenMined/TenSEAL}

\bibitem{bolt2024}
Q.~Pang, J.~Zhu, H.~M\"{o}llering, W.~Zheng, and T.~Schneider,
``BOLT: Privacy-Preserving, Accurate and Efficient Inference for
Transformers,''
in \emph{IEEE S\&P 2024}, pp.~133--133.

\bibitem{bumblebee2025}
W.-J.~Lu, Z.~Huang, Z.~Gu, J.~Li, J.~Liu, K.~Ren, C.~Hong, T.~Wei,
and W.~Chen,
``Bumblebee: Secure Two-Party Inference Framework for Large Transformers,''
in \emph{NDSS 2025}.

\bibitem{bonawitz2017}
K.~Bonawitz et al.,
``Practical Secure Aggregation for Privacy-Preserving Machine Learning,''
in \emph{Proc. ACM CCS 2017}, pp.~1175--1191.

\bibitem{hear2022}
J.~Ma et al.,
``HEAR: Human Action Recognition via Neural Networks on Homomorphically
Encrypted Data,''
\emph{arXiv preprint arXiv:2104.09164}, 2022.

\bibitem{chengu2024}
G.~Chen and Q.~Wu,
``A Review of Personalized Federated Reinforcement Learning,''
\emph{Int.\ J.\ Comput.\ Sci.\ Inf.\ Technol.}, vol.~3, no.~1,
pp.~1--9, 2024.

\bibitem{cicids2017}
I.~Sharafaldin, A.~H.~Lashkari, and A.~A.~Ghorbani,
``Toward Generating a New Intrusion Detection Dataset and Intrusion Traffic
Characterization,''
in \emph{Proc. ICISSP 2018}, pp.~108--116.

\bibitem{HEstandard}
M.~Albrecht et al.,
``Homomorphic Encryption Security Standard,''
HomomorphicEncryption.org, Tech. Rep., 2021.
\url{https://homomorphicencryption.org/standard}

\end{thebibliography}

\end{document}
\typeout{get arXiv to do 4 passes: Label(s) may have changed. Rerun}